% Intended LaTeX compiler: pdflatex
\documentclass[a4paper,12pt]{article}
\usepackage[utf8]{inputenc}
\usepackage[T1]{fontenc}
\usepackage{graphicx}
\usepackage{longtable}
\usepackage{wrapfig}
\usepackage{rotating}
\usepackage[normalem]{ulem}
\usepackage{amsmath}
\usepackage{amssymb}
\usepackage{capt-of}
\usepackage{hyperref}
\usepackage[margin=0.2in]{geometry}
\usepackage{amsmath}
\usepackage{amsfonts}
\usepackage{indentfirst}
\usepackage{pgffor}
\usepackage{amssymb}
\usepackage{pifont}
\usepackage{cancel}
\usepackage{amsthm}
\usepackage{polynom}
\usepackage{tikz}
\usepackage[tikz]{bclogo}
\usepackage{listings}
\usepackage[most]{tcolorbox}
\usepackage{forest}
\usepackage{adjustbox}
\usepackage{tikz-3dplot}
\usepackage{mathtools}
\usepackage{pgfplots}
\usetikzlibrary{tikzmark,calc,fit,matrix,arrows,automata,positioning}
\usepackage{centernot}
\usepackage{lmodern}
\usepackage{physics}
\usepackage[boxed,linesnumbered,vlined]{algorithm2e}
\usepackage{algpseudocode}
\usepackage{algorithmicx}
\usepackage{svg}
\tikzstyle{every picture}+=[remember picture]
\DeclareMathOperator{\spn}{span}
\DeclareMathOperator{\dom}{domain}
\DeclareMathOperator{\ran}{range}
\DeclareMathOperator{\rng}{range}
\DeclareMathOperator{\img}{Im}
\DeclareMathOperator{\adj}{adj}
\DeclareMathOperator{\Var}{Var}
\DeclareMathOperator{\cov}{cov}
\newcommand\dif[1]{\:\textrm{d}#1}
\DeclarePairedDelimiter\ceil{\lceil}{\rceil}
\DeclarePairedDelimiter\floor{\lfloor}{\rfloor}
\newcommand{\ihat}{\hat{\textbf{\i}}}
\newcommand{\jhat}{\hat{\textbf{\j}}}
\newcommand{\khat}{\hat{\textbf{k}}}
\newcommand{\rhat}{\hat{\textbf{r}}}
\newcommand{\thetahat}{\boldsymbol{\hat{\theta}}}
\newcommand{\?}{\stackrel{?}{=}}
\newcommand{\cmark}{\text{\ding{51}}}
\newcommand{\xmark}{\text{\ding{55}}}
\author{mahmood}
\date{\today}
\title{code tree}
\hypersetup{
 pdfauthor={mahmood},
 pdftitle={code tree},
 pdfkeywords={},
 pdfsubject={},
 pdfcreator={Emacs 30.0.50 (Org mode 9.6)}, 
 pdflang={English}}
\begin{document}

\maketitle
\definecolor{Brown}{cmyk}{0,0.81,1,0.60}
\definecolor{OliveGreen}{cmyk}{0.64,0,0.95,0.40}
\definecolor{CadetBlue}{cmyk}{0.62,0.57,0.23,0}
\definecolor{lightlightgray}{gray}{0.9}
\lstset{
  language=C,                             % Code langugage
  basicstyle=\ttfamily,                   % Code font, Examples: \footnotesize, \ttfamily
  keywordstyle=\color{OliveGreen},        % Keywords font ('*' = uppercase)
  commentstyle=\color{gray},              % Comments font
  numbers=left,                           % Line nums position
  numberstyle=\tiny,                      % Line-numbers fonts
  stepnumber=1,                           % Step between two line-numbers
  numbersep=5pt,                          % How far are line-numbers from code
  backgroundcolor=\color{lightlightgray}, % Choose background color
  frame=single,                           % A frame around the code
  tabsize=2,                              % Default tab size
  captionpos=b,                           % Caption-position = bottom
  breaklines=true,                        % Automatic line breaking?
  breakatwhitespace=false,                % Automatic breaks only at whitespace?
  showspaces=false,                       % Dont make spaces visible
  showtabs=false,                         % Dont make tabls visible
  columns=flexible,                       % Column format
  morekeywords={__global__, __device__},  % CUDA specific keywords
  showstringspaces=false,                 % dont show spaces in strings
}
% so that latex doesnt complain about sage not being a language
\lstdefinelanguage{sage}{}

% environment definition for special blocks
\theoremstyle{definition}
\newtheorem{my_example}{Example}
\counterwithin{my_example}{section}
\counterwithin{my_example}{subsection}
\newtheorem{definition}{Definition}
\counterwithin{definition}{section}
\counterwithin{definition}{subsection}
\newtheorem{characteristic}{Characteristic}
\newtheorem{entailment}{Entailment}
%\newtheore*{proof}{Proof}
\newtheorem{lemma}{Lemma}
\newtheorem{question}{Question}
\newtheorem{subquestion}{Subquestion}[question]
%\newtheorem{note}{Note}
\newtheorem{answer}{Answer}
\counterwithin{answer}{question}
\counterwithin{answer}{subquestion}
\newtheorem{step}{Step}[answer]

% the following snippet auto resizes images to fit properly
% Determine if the image is too wide for the page.
% \makeatletter
% \def\ScaleIfNeeded{%
%   \ifdim\Gin@nat@width>\linewidth
%     \linewidth
%   \else
%     \Gin@nat@width
%   \fi
% }
% \makeatother
% % Resize figures that are too wide for the page.
% \let\oldincludegraphics\includegraphics
% \renewcommand\includegraphics[2][]{%
%   \oldincludegraphics[width=\ScaleIfNeeded]{#2}
% }

\newenvironment{note}
  {\begin{bclogo}[logo=\bcinfo, couleurBarre=orange, couleur=lightgray]{ note}}
  {\end{bclogo}}
\newenvironment{attention}
  {\begin{bclogo}[logo=\bcattention, couleurBarre=red, noborder=true, 
                couleur=red!15]{Important!}}
  {\end{bclogo}}

let \(k\) be a \href{20230119174917-prefix_code.org}{prefix code} over the characters \(C=\{c_1,\dots,c_n\}\), the \textbf{code tree} of \(k\), \(T_k=(V_k,E_k)\) is a \href{data_structures.org}{binary tree} such that for every inner node \(v \in V_k\), one of the lines coming out of \(v\) are labelled by 0, and the other by 1. if from \(v\) comes out one line only, then its labelled either 0 or 1.\\[0pt]
each character in \(C\) is matched with a unique root-to-leaf path in \(T_k\), we denote the path that corresponds to \(c_i\) by \(a_i\). the leaf at the end of the path \(a_i\) is tagged by \(c_i\) and the lines through \(a_i\), starting from the root to the leaf are tagged by the bits of the code \(k(c_i)\)\\[0pt]
in other words: if \(k(c_i)=b_1,\dots,b_l\) then in \(a_i\) we have \(l\) lines, and the line of \(j\), \(1 \le j \le l\), is tagged by \(b_j\). the code trees of \(k_1\) and of \(k_2,T_{k_1},T_{k_2}\) are shown in the following diagrams\\[0pt]
\begin{my_example}
k\textsubscript{1}:\\[0pt]
\includesvg{/Users/mahmooz/brain/out/oMu1J1N}\\[0pt]

\begin{tabular}{lrr}
char & freq & \(k_1\) \href{20230119170007-fixed_length_code_word.org}{fixed-length code word}\\[0pt]
\hline
a & 45 & 000\\[0pt]
b & 13 & 001\\[0pt]
c & 12 & 010\\[0pt]
d & 16 & 011\\[0pt]
e & 9 & 100\\[0pt]
f & 5 & 101\\[0pt]
\end{tabular}

k\textsubscript{2}:\\[0pt]

\includesvg{/Users/mahmooz/brain/out/Vm959rL}\\[0pt]

\begin{tabular}{lrrr}
char & freq & \(k_1\), fixed-length & \(k_2\), variable-length\\[0pt]
\hline
a & 45 & 000 & 0\\[0pt]
b & 13 & 001 & 101\\[0pt]
c & 12 & 010 & 100\\[0pt]
d & 16 & 011 & 111\\[0pt]
e & 9 & 100 & 1101\\[0pt]
f & 5 & 101 & 1100\\[0pt]
\end{tabular}

a tree code \(T_k\) allows simple and efficient decoding of a file that has been encoded using \(k\), such that each character in the original file is decoded by processing \(T_k\) from the root to one of the leaves\\[0pt]
for example, if we were to look in the trees we'd find that \(k_1(adf)=000\circ011\circ101=000011101\), to decode the first character in \(k_1(adf)\), we process the tree starting from the root \(T_{k_1}\), we progress using the lines that are tagged by 0,0,0, which are the 3 first bits in the encoded file, we arrive to the leaf that is tagged by \(a\), which is the first character in the encoded file\\[0pt]
now the first 3 bits are removed from the encoded file and as such the decoding of the first character is done. continuing, the decoding of each of the characters is done by processing a root-to-leaf path, addition of the character that appears in the leaf to the decoded file, and removal of the bits from the original encoded file\\[0pt]
\end{my_example}
\section{cost of code tree}
\label{sec:orgd7d6a6d}
let \(F\) be a file over the set of characters \(C=\{c_1,c_2,\dots,c_n\}\), let \(k\) be \href{20230119174917-prefix_code.org}{prefix code} of \(C\), with the code tree \(T_k\), for each \(c_i \in C\), we denote by \(f_i\) the frequency of \(c_i\) in \(F\), and recall that the number of bits in \(k(c_i)\) is equal to the number of lines of \(a_i\). the \textbf{cost} of \(T_k\), in correspondance to the file \(F\), \(\mathrm{cost}(T_k,F)\), is defined as the number of bits in \(k(F)\), so:\\[0pt]
\[
\mathrm{cost}(T_k,F) = \sum_{i=1}^{n} f_i \cdot |k(c_i)| = \sum_{i=1}^{n} f_i \cdot |a_i| 
\]\\[0pt]
the code tree \(T_k\) is optimal in correspondance to \(F\), if there exists no \href{20230119174917-prefix_code.org}{prefix code} \(k'\) such that \(\mathrm{cost}(T_k,F) < \mathrm{cost}(T_k,F)\)\\[0pt]
\section{representative leaf and leaf array}
\label{sec:org9e58d4f}
let \(F\) be a file over the set of characters \(C=\{c_1,c_2,\dots,c_n\}\), for each \(1 \le i \le n\), let \(f_i\) be the frequency of \(c_i\). the leaf that \textbf{represents} the character \(c_i\) is the node labelled by the \href{discrete_maths2.org}{ordered pair} \(p_i=(c_i,f_i)\). the components of \(p_i\) are denoted by \(p_i.f\) and \(p_i.c\)\\[0pt]
the \textbf{leaf array} of \(F\) is an array \(D[]\) of \(n\) elements, such that \(D[i]\) stores the leaf that represents \(c_i\), i.e. \((\forall 1 \le i \le n)[D[i]=p_i]\)\\[0pt]
\begin{note}
the leaf array contains all the necessary data to calculate the optimal code tree for \(F\)\\[0pt]
\end{note}
\end{document}